\documentclass[12pt]{article}
\usepackage[utf8]{inputenc}
\usepackage[T1]{fontenc}
\usepackage[margin=1in]{geometry}
\usepackage{setspace}
\usepackage{xcolor}
\usepackage[colorlinks=true, urlcolor=blue!60!black, linkcolor=blue!60!black]{hyperref}
\setlength{\parindent}{2em}
\setlength{\parskip}{1.25em}
\renewcommand{\baselinestretch}{1.0}
\begin{document}
\begin{center}
{\LARGE \textit{South-Central Timber Development, Inc. v. Wunnicke}, 467 U.S. 82 (1984)\par}
\vspace{0.8em}
{\large Case Brief\par}
\end{center}
\vspace{1em}

\section{Memory Jogger}
Alaska offered state timber for sale on condition that buyers process it in Alaska before export; the Court held that a parallel federal-lands policy was not ``unmistakably clear'' congressional authorization, and a plurality held that the market-participant exception does not let a State impose conditions that regulate a downstream market it does not participate in (\href{https://supreme.justia.com/cases/federal/us/467/82/}{Justia}).

\section{Detailed Case Facts}
In September 1980, the Alaska Department of Natural Resources published notice of a planned October 23, 1980 sale of approximately 49 million board-feet of state-owned standing timber near Icy Cape, Alaska (\href{https://supreme.justia.com/cases/federal/us/467/82/}{467 U.S.\ at 84}). Pursuant to 11 Alaska Admin.\ Code \S\ 76.130 (1974) (recodified at \S\S\ 71.230, 71.910 (1982)) and a May 1974 Governor's Policy Statement, the notice of sale, prospectus, and proposed contract (Section~68) required ``primary manufacture'' of the timber within Alaska before the winning bidder could ship it out of state (\href{https://supreme.justia.com/cases/federal/us/467/82/}{467 U.S.\ at 84--85 \& nn.\ 1--2}).

The condition applied solely by contract to state-owned timber, not private timber, and sought to protect existing industries, foster new industries, derive revenue, and manage forests on a sustained-yield basis (\href{https://supreme.justia.com/cases/federal/us/467/82/}{467 U.S.\ at 85}). Because the restriction depressed bids, Alaska charged a significantly lower timber price than it otherwise would have obtained. Buyers typically complied by sawing logs into ``cants'' (slabbed to a maximum thickness of 12 inches, later $8\frac{3}{4}$ inches, or squared on four sides), which wasted valuable wood and made cants worth less than round logs and virtually unsaleable within Alaska (\href{https://supreme.justia.com/cases/federal/us/467/82/}{467 U.S.\ at 85--86 \& nn.\ 3, 5}).

Petitioner South-Central Timber Development, Inc., an Alaska corporation without an in-state sawmill, logged standing timber and exported unprocessed round logs almost exclusively to Japan, which absorbed over 90\% of Alaska timber due to high domestic shipping costs (\href{https://supreme.justia.com/cases/federal/us/467/82/}{467 U.S.\ at 85--86 \& n.\ 4}). Separately, since 1928, federal Forest Service regulations (36 C.F.R.\ \S\ 223.10(c) (1983)) and congressional appropriations riders had restricted exports of unprocessed timber cut from \emph{federal} lands in Alaska and western States (\href{https://supreme.justia.com/cases/federal/us/467/82/}{467 U.S.\ at 88--89}).

\section{Procedural History}
South-Central sued in the U.S.\ District Court for the District of Alaska to enjoin enforcement of the primary-manufacture clause under the Dormant Commerce Clause. The District Court agreed and issued an injunction (\textit{South-Central Timber Development, Inc.\ v.\ LeResche}, 511 F.\ Supp.\ 139 (D.\ Alaska 1981)) (\href{https://supreme.justia.com/cases/federal/us/467/82/}{467 U.S.\ at 86--87}). The Ninth Circuit reversed (693 F.2d 890 (9th Cir.\ 1982)), declining to reach the market-participant or substantive Commerce Clause issues because it concluded that federal primary-manufacture policies on federal lands ``implicitly authorized'' Alaska's parallel policy for state lands (\href{https://supreme.justia.com/cases/federal/us/467/82/}{467 U.S.\ at 87--88}). The Supreme Court granted certiorari (464 U.S.\ 890 (1983)), with the United States appearing as \textit{amicus curiae} supporting South-Central (\href{https://supreme.justia.com/cases/federal/us/467/82/}{467 U.S.\ at 84, 90 n.\ 6}).

\section{Judicial Votes}
The Supreme Court \textbf{reversed and remanded} the Ninth Circuit by a \textbf{6--2} vote, with one Justice not participating. The case was argued on February 29, 1984, and decided on May 22, 1984 (\href{https://supreme.justia.com/cases/federal/us/467/82/}{Justia}; \href{https://www.law.cornell.edu/supremecourt/text/467/82}{Cornell LII}):

\begin{itemize}
        \item \textbf{Majority Opinion (Parts I and II --- Congressional Authorization):} Justice Byron R.\ White delivered the opinion of the Court, joined by Chief Justice Warren E.\ Burger and Justices William J.\ Brennan, Jr., Harry A.\ Blackmun, Lewis F.\ Powell, Jr., and John Paul Stevens (6 Justices).
        \item \textbf{Plurality Opinion (Parts III and IV --- Market-Participant Exception \& Commerce Clause Merits):} Justice White announced the judgment of the Court and delivered a plurality opinion, joined by Justices Brennan, Blackmun, and Stevens (4 Justices).
        \item \textbf{Concurrence:} Justice Brennan filed a concurring opinion (467 U.S.\ at 101).
        \item \textbf{Concurrence in Part and Concurrence in the Judgment:} Justice Powell filed an opinion concurring in part and concurring in the judgment, joined by Chief Justice Burger (467 U.S.\ at 101).
        \item \textbf{Dissent:} Justice William H.\ Rehnquist filed a dissenting opinion, joined by Justice Sandra Day O'Connor (467 U.S.\ at 101--103).
        \item \textbf{Not Participating:} Justice Thurgood Marshall took no part in the decision of the case.

\end{itemize}

\section{Holding}
\textbf{Questions Presented:}

\begin{enumerate}
    \item \textbf{Congressional Authorization (Majority):} Did federal statutes and regulations requiring in-state primary manufacture of timber harvested from \emph{federal} lands in Alaska implicitly authorize Alaska to impose a parallel local-processing mandate on \emph{state}-owned timber?
    \item \textbf{Market-Participant Exception --- Downstream Limit (Plurality):} Absent congressional authorization, was Alaska exempt from Dormant Commerce Clause scrutiny as a ``market participant'' when it used its sale of state-owned timber to require purchasers to process the timber in-state before export?
    \item \textbf{Dormant and Foreign Commerce Clause Merits (Plurality):} If the market-participant exception does not apply, does Alaska's in-state processing requirement violate the Dormant Commerce Clause and heightened Foreign Commerce Clause scrutiny?
\end{enumerate}

\textbf{Holding:}

\begin{enumerate}
    \item \textbf{No Implicit Congressional Authorization (6--2 Majority, Parts I--II):} \textbf{No.} To remove a state restraint from the Dormant Commerce Clause, congressional intent must be ``unmistakably clear.'' A parallel federal policy governing federal lands does not authorize a State to impose a comparable restraint on state lands (\href{https://supreme.justia.com/cases/federal/us/467/82/}{467 U.S.\ at 91--93}).
    \item \textbf{Downstream Limit on Market-Participant Immunity (4-Justice Plurality, Part III):} \textbf{No.} The market-participant doctrine permits a State to burden commerce only within the narrowly defined market in which it participates; it does not immunize contractual, statutory, or regulatory conditions exerting a substantial regulatory effect \emph{downstream} in a separate market (timber processing) in which the State does not participate (\href{https://supreme.justia.com/cases/federal/us/467/82/}{467 U.S.\ at 97--99}).
    \item \textbf{Virtual Per Se Invalidity (4-Justice Plurality, Part IV):} \textbf{Yes.} Alaska's naked local-processing mandate is protectionist and virtually \textit{per se} invalid under the Dormant Commerce Clause, reinforced by heightened scrutiny for burdens on foreign commerce (\href{https://supreme.justia.com/cases/federal/us/467/82/}{467 U.S.\ at 99--101}).
\end{enumerate}

\section{Analysis of Opinions}
\subsection{Opinion of the Court and Plurality Opinion (Justice White)}

\subsubsection{Part II (Opinion of the Court): ``Unmistakably Clear'' Congressional Authorization}

While Congress may redefine the distribution of power over commerce and permit otherwise invalid state restraints, precedents such as \href{https://supreme.justia.com/cases/federal/us/458/941/}{\textit{Sporhase v.\ Nebraska ex rel.\ Douglas}}, 458 U.S.\ 941, 960 (1982), and \href{https://supreme.justia.com/cases/federal/us/455/331/}{\textit{New England Power Co.\ v.\ New Hampshire}}, 455 U.S.\ 331, 341 (1982), require affirmative congressional contemplation (\href{https://supreme.justia.com/cases/federal/us/467/82/}{467 U.S.\ at 87--91}). Rejecting Alaska's reading of \href{https://supreme.justia.com/cases/federal/us/460/204/}{\textit{White v.\ Massachusetts Council of Construction Employers, Inc.}}, 460 U.S.\ 204 (1983), Justice White explained that although ``expressly stated'' is not talismanic, congressional intent must be ``unmistakably clear.'' This clear-statement rule protects unrepresented out-of-state interests against economic Balkanization (\href{https://supreme.justia.com/cases/federal/us/441/322/}{\textit{Hughes v.\ Oklahoma}}, 441 U.S.\ 322, 325 (1979)) by ensuring any authorization is a collective national decision in Congress---a need heightened where state policy affects foreign relations (\href{https://supreme.justia.com/cases/federal/us/467/82/}{467 U.S.\ at 91--92 \& n.\ 7}). Because Congress legislated only for federal lands (and had declined to authorize state processing mandates, \href{https://supreme.justia.com/cases/federal/us/467/82/}{467 U.S.\ at 93 n.\ 8}), harmony with federal goals was insufficient.

\subsubsection{Part III (Plurality Opinion): The Downstream Market-Participant Boundary}

The plurality distinguished the Court's three prior market-participant decisions:
\begin{itemize}
    \item \href{https://supreme.justia.com/cases/federal/us/426/794/}{\textit{Hughes v.\ Alexandria Scrap Corp.}}, 426 U.S.\ 794 (1976): Maryland paid direct scrap bounties, leaving sellers free to forgo the subsidy and export; Alaska stripped buyers of that choice (\href{https://supreme.justia.com/cases/federal/us/467/82/}{467 U.S.\ at 93--95}).
    \item \href{https://supreme.justia.com/cases/federal/us/447/429/}{\textit{Reeves, Inc.\ v.\ Stake}}, 447 U.S.\ 429 (1980): South Dakota's cement preference lacked all three limiting features present here---foreign commerce, a natural resource, and post-sale resale/disposition restrictions (\href{https://supreme.justia.com/cases/federal/us/467/82/}{467 U.S.\ at 95--96}).
    \item \href{https://supreme.justia.com/cases/federal/us/460/204/}{\textit{White v.\ Massachusetts Council of Construction Employers, Inc.}}, 460 U.S.\ 204 (1983): Boston's resident-hiring rule reached subcontractors' workers only because everyone affected was informally ``working for the city'' during an ongoing public project (\href{https://supreme.justia.com/cases/federal/us/465/208/}{\textit{United Building \& Construction Trades Council v.\ Mayor of Camden}}, 465 U.S.\ 208, 219 (1984)), whereas Alaska forced buyers after delivery and payment to deal with a ``stranger to the contract'' (\href{https://supreme.justia.com/cases/federal/us/467/82/}{467 U.S.\ at 94--97, 99}).
\end{itemize}
Defining the market narrowly (timber sales, not timber processing), the plurality held that a State cannot leverage proprietary participation in one market to regulate downstream private relationships in another. Analogizing to common-law restraints on alienation and antitrust limits on vertical restraints, Justice White rejected Alaska's ``equivalent means'' defense (that it could vertically integrate, pay direct subsidies, or sell only to local processors) and barred formalistic evasion via title-retention clauses (\href{https://supreme.justia.com/cases/federal/us/467/82/}{467 U.S.\ at 97--99 \& n.\ 11}).

\subsubsection{Part IV (Plurality Opinion): Per Se Invalidity and Foreign Commerce}

As a naked local-processing mandate designed to promote local employment, Alaska's rule fell within the virtual \textit{per se} invalidity rule of \href{https://supreme.justia.com/cases/federal/us/397/137/}{\textit{Pike v.\ Bruce Church, Inc.}}, 397 U.S.\ 137, 145 (1970), \href{https://supreme.justia.com/cases/federal/us/278/1/}{\textit{Foster-Fountain Packing Co.\ v.\ Haydel}}, 278 U.S.\ 1 (1928), and \href{https://supreme.justia.com/cases/federal/us/437/617/}{\textit{City of Philadelphia v.\ New Jersey}}, 437 U.S.\ 617, 624 (1978). That conclusion was reinforced by heightened Foreign Commerce Clause scrutiny, because the Federal Government must ``speak with one voice'' in international trade (\href{https://supreme.justia.com/cases/federal/us/441/434/}{\textit{Japan Line, Ltd.\ v.\ County of Los Angeles}}, 441 U.S.\ 434 (1979)) (\href{https://supreme.justia.com/cases/federal/us/467/82/}{467 U.S.\ at 99--101}).

\subsection{Concurring Opinion (Justice Brennan)}

Justice Brennan joined Justice White in full while observing that the clash between the plurality and dissent exposed ``the inherent weakness'' of the market-participant doctrine (citing his dissent in \href{https://supreme.justia.com/cases/federal/us/426/794/}{\textit{Alexandria Scrap}}, 426 U.S.\ at 817).

\subsection{Opinion Concurring in Part and in the Judgment (Justice Powell)}

Justice Powell, joined by Chief Justice Burger, joined Parts~I and II but would have remanded to the Ninth Circuit to decide the market-participant and \href{https://supreme.justia.com/cases/federal/us/397/137/}{\textit{Pike}} issues in the first instance.

\subsection{Dissenting Opinion (Justice Rehnquist)}

Justice Rehnquist, joined by Justice O'Connor, criticized the plurality's market-participant/regulator line as ``both artificial and unconvincing.'' Noted the irony of citing antitrust cases (which apply to States only as market participants, whereas regulators enjoy \textit{Parker v.\ Brown} immunity) to label Alaska a regulator, dismissed the downstream test as circular, and argued Alaska was merely paying buyers via a lower price to hire local processors---a result functionally identical to permissible alternatives under \href{https://supreme.justia.com/cases/federal/us/447/429/}{\textit{Reeves}}, \href{https://supreme.justia.com/cases/federal/us/426/794/}{\textit{Alexandria Scrap}}, and \href{https://supreme.justia.com/cases/federal/us/460/204/}{\textit{White}}.

\subsection{Subsequent Developments}

\begin{itemize}

\item \textbf{The ``Unmistakably Clear'' Canon Applied:} First, Part~II's six-Justice majority cemented the ``unmistakably clear'' canon for congressional authorization of state trade barriers, relied upon in \href{https://supreme.justia.com/cases/federal/us/477/131/}{\textit{Maine v.\ Taylor}}, 477 U.S.\ 131, 138--139 (1986) (rejecting implied consent under the Lacey Act), and contrasted with express statutory authorization in \href{https://supreme.justia.com/cases/federal/us/472/159/}{\textit{Northeast Bancorp, Inc.\ v.\ Board of Governors}}, 472 U.S.\ 159 (1985).

\item \textbf{The Downstream Limit in Later Market-Participant Cases:} Second, although Part~III commanded only a four-Justice plurality, courts treat its ``downstream regulation'' limit as the canonical outer boundary of the market-participant exception. Subsequent decisions refused to extend market-participant immunity to discriminatory taxes or exemptions (\href{https://supreme.justia.com/cases/federal/us/486/269/}{\textit{New Energy Co.\ of Indiana v.\ Limbach}}, 486 U.S.\ 269 (1988); \href{https://supreme.justia.com/cases/federal/us/520/564/}{\textit{Camps Newfound/Owatonna, Inc.\ v.\ Town of Harrison}}, 520 U.S.\ 564 (1997)), while carving out a separate safe harbor only for laws favoring \emph{public} entities over all private competitors alike (\href{https://supreme.justia.com/cases/federal/us/550/330/}{\textit{United Haulers Ass'n v.\ Oneida-Herkimer Solid Waste Management Authority}}, 550 U.S.\ 330 (2007); \href{https://supreme.justia.com/cases/federal/us/553/328/}{\textit{Department of Revenue of Kentucky v.\ Davis}}, 553 U.S.\ 328 (2008)).

\end{itemize}

\section{Examples: Future Applications}
\begin{enumerate}

\item \textbf{Same-Side (Invalid Downstream Mandate): In-State Refining of State-Land Lithium.} Nevada auctions state-land lithium leases requiring winning bidders to refine all extracted lithium at in-state facilities before export, invoking a federal statute requiring domestic processing of \emph{federal}-land minerals. Under \textit{Wunnicke} Part~II, parallel federal-land policy is not ``unmistakably clear'' authorization. Under Part~III, Nevada participates in raw-lithium leasing, not downstream refining; imposing post-extraction processing strings on a natural resource entering foreign trade fails market-participant immunity and is virtually \textit{per se} invalid under \href{https://supreme.justia.com/cases/federal/us/397/137/}{\textit{Pike v.\ Bruce Church, Inc.}}, 397 U.S.\ 137 (1970).

\item \textbf{Same-Side (Invalid Downstream Mandate): Resale Restrictions on State-Made Cement.} A State that operates its own cement plant sells cement to private distributors on condition that they resell it only to customers within the State, even after title and payment have passed. Out-of-state builders who would buy from the distributors challenge the condition. The State may choose its own customers, as \textit{Reeves} allowed, but here it reaches past its sale to control the buyers' later dealings in a resale market where the State is not a participant. The plurality identified exactly this kind of post-sale restriction as one of the features absent from \textit{Reeves}, and a resale ban that walls off the State's market for local benefit is the kind of downstream regulation \textit{Wunnicke} forbids.

\item \textbf{Opposite-Side (Valid Market Participation): A State-Owned Refinery Choosing Its Customers.} Return to the lithium hypothetical, but suppose Nevada instead operates a state-owned refinery and sells refined lithium exclusively to Nevada battery makers during shortages. Nevada is now acting as a seller in the very market it has entered and is simply choosing its customers, the conduct \textit{Reeves} protected when South Dakota sold cement from its own plant to residents first. No condition follows the refined lithium after sale. If Nevada buys raw lithium on the open market and preferentially sells only its manufactured output, the preference likely falls within the market-participant exception. A preference that effectively reserves access to state-owned raw lithium would raise the natural-resources issue \textit{Reeves} reserved.

\item \textbf{Opposite-Side (Valid Market Participation): Cash Subsidies Instead of Export Conditions.} Alternatively, suppose Nevada auctions raw leases without export conditions while paying direct general-fund cash subsidies to in-state refiners. Lessees remain free to ship raw ore anywhere, so the State has not stripped buyers of their choice, the feature that distinguished \textit{Alexandria Scrap} from Alaska's contract clause. The subsidy may favor local refiners, but a direct payment from the State's own funds is not a regulation of the downstream market, and the plurality noted that such alternatives remain open to a State.

\item \textbf{Fence-Sitter (Close Call): Requiring Resident Logging Crews on State Land.} Alaska sells standing timber on state land and requires that the buyer use crews made up of Alaska residents to cut and haul the timber off state land; once the logs leave state land, the buyer may process and ship them anywhere. Alaska argues that, as in \textit{White}, the workers are effectively working on a state project before the sale is complete. The challenger replies that the condition regulates the separate labor market for logging services. Because the plurality defined the market by intuition without an objective test, and because Part~III never commanded a majority, it is unclear whether harvesting on state land falls inside the timber-sales market or downstream of it.

\end{enumerate}

\section{Critique}
\textbf{Market-Definition Indeterminacy:} The plurality defines the ``relevant market'' narrowly (raw timber vs.\ processing) by ``intuition'' (\href{https://supreme.justia.com/cases/federal/us/467/82/}{467 U.S.\ at 98}) without an objective economic test.

\textbf{Formalism vs.\ Political Transparency:} Justice Rehnquist argued Alaska could achieve the same protectionism via vertical integration or direct subsidies (\href{https://supreme.justia.com/cases/federal/us/467/82/}{467 U.S.\ at 103}). Yet \textit{Wunnicke}'s line forces visible legislative appropriations rather than hidden resource-auction discounts.

\textbf{Plurality Fragility:} Only four Justices joined Parts~III--IV; Justice Powell and Chief Justice Burger joined only Parts~I--II and the judgment, leaving Part~III formally non-majority precedent.

\section{Key Quotations}
\begin{quote}
``There is no talismanic significance to the phrase `expressly stated,' however; it merely states one way of meeting the requirement that, for a state regulation to be removed from the reach of the dormant Commerce Clause, congressional intent must be unmistakably clear. The requirement that Congress affirmatively contemplate otherwise invalid state legislation is mandated by the policies underlying dormant Commerce Clause doctrine.''
\end{quote}
\noindent --- Justice White, \href{https://supreme.justia.com/cases/federal/us/467/82/}{467 U.S.\ at 91--92} (\href{https://www.law.cornell.edu/supremecourt/text/467/82}{Cornell LII}).

\begin{quote}
``The fact that the state policy in this case appears to be consistent with federal policy---or even that state policy furthers the goals we might believe that Congress had in mind---is an insufficient indicium of congressional intent. Congress acted only with respect to federal lands; we cannot infer from that fact that it intended to authorize a similar policy with respect to state lands.''
\end{quote}
\noindent --- Justice White, \href{https://supreme.justia.com/cases/federal/us/467/82/}{467 U.S.\ at 92--93} (\href{https://www.law.cornell.edu/supremecourt/text/467/82}{Cornell LII}).

\begin{quote}
``The limit of the market participant doctrine must be that it allows a State to impose burdens on commerce within the market in which it is a participant, but allows it to go no further. The State may not impose conditions, whether by statute, regulation, or contract, that have a substantial regulatory effect outside of that particular market. Unless the `market' is relatively narrowly defined, the doctrine has the potential of swallowing up the rule that States may not impose substantial burdens on interstate commerce even if they act with the permissible state purpose of fostering local industry.''
\end{quote}
\noindent --- Justice White (plurality opinion), \href{https://supreme.justia.com/cases/federal/us/467/82/}{467 U.S.\ at 97--98} (\href{https://www.law.cornell.edu/supremecourt/text/467/82}{Cornell LII}).

\begin{quote}
``Alaska is merely paying the buyer of the timber indirectly, by means of a reduced price, to hire Alaska residents to process the timber. Under existing precedent, the State could accomplish that same result in any number of ways. \dots\ It seems to me unduly formalistic to conclude that the one path chosen by the State as best suited to promote its concerns is the path forbidden it by the Commerce Clause.''
\end{quote}
\noindent --- Justice Rehnquist (dissenting), \href{https://supreme.justia.com/cases/federal/us/467/82/}{467 U.S.\ at 103} (\href{https://www.law.cornell.edu/supremecourt/text/467/82}{Cornell LII}).


\end{document}
